% introduction.tex -- Introduction
\section{Introduction}

During the October 1987 market decline, portfolio insurance concentrated
selling across participating traders~\cite{shiller1988portfolio,grossman1988analysis}.
The forced unwind of leveraged funds transmitted losses across
quantitative equity strategies in August 2007~\cite{khandani2011what}.
Such episodes motivate a common modeling question: how can a change in
the behavior of some traders become a market-wide event? Agent-based
models address collective behavior through imitation, strategy adaptation,
and responses to a shared price~\cite{lebaron2006agent,hommes2006heterogeneous,samanidou2007agent,axtell2025agent}.
Herding models describe how individual decisions produce collective
activity~\cite{kirman1993ants,cont2000herd,lux1999scaling}, while lattice
models restrict influence to neighboring traders~\cite{bornholdt2001ising,bartolozzi2004stochastic,qiu2007cellular}.
These interactions can generate large fluctuations and clustered
volatility, two features of financial returns~\cite{cont2001empirical}.
However, reproducing aggregate fluctuations does not by itself explain
how an established pattern of trading responds when part of the
population changes its behavior. The size of the changed group matters,
but so does the nature of the change. A temporary increase in selling
and a persistent change in how traders respond to their neighbors need
not produce the same collective outcome.

A cellular automaton makes that distinction concrete, because it gives
each trader an explicit local decision rule. Traders occupy Buy, Hold, or
Sell states on a square lattice and update from their four nearest
neighbors, and a totalistic rule assigns the next state from the sum of
those neighbor states~\cite{wolfram1984universality,packard1985two,wolfram2002new}.
Local interactions of this kind have been used to generate market
fluctuations and regime changes~\cite{wei2003cellular,bartolozzi2004stochastic,fan2009effect,atman2012influence}.
The lookup table makes a persistent behavioral change precise: one
response can change while the remaining entries stay fixed. A one-time
Sell pulse changes only the current configuration, whereas changing the
rule changes the response whenever the affected neighborhood recurs, and
the altered response remains in place even after a trader leaves the
Sell state. Following matched trajectories from the same realized
lattice allows the effect of that persistent change to be measured
against continued baseline
behavior~\cite{glasserman1992guidelines,law2015simulation}.
The question is then how many traders must adopt the changed rule
before local selling grows into a Sell majority, and whether that
majority can develop without feedback from the market price.

Whether a Sell majority forms is one question. What it does to the price
is a second question, and the answer depends on the resources available
to execute the resulting orders. A trader needs shares to sell or cash
to buy, and an eligible order needs a counterparty. Dealer inventory and
capital therefore constrain the unmatched flow that can
execute~\cite{garman1976market,stoll1978supply,ho1981optimal,madhavan2000market,adrian2017dealer}.
Price-impact models connect that executed flow to changes in the
quote~\cite{kyle1985continuous,hasbrouck1991measuring,farmer2005predictive,bouchaud2009how}.
Agent-based markets have incorporated finite endowments, order books,
leverage limits, and dealers with constrained balance
sheets~\cite{raberto2001agent,chiarella2002simulation,thurner2012leverage,braunmunzinger2018agent}.
For a lattice market, separating desired trades from executed orders
allows a Sell-dominated state to coexist with a limited price decline.
The execution constraint also matters when traders respond to the price,
because restricting unmatched flow changes the quote and therefore
changes the signal used in subsequent decisions. Dealer capacity can
then influence both the spread of selling and its financial
consequences. An unlimited dealer can absorb all eligible unmatched
flow, whereas a finite dealer may reach its cash or inventory limit
while many traders remain in Sell.

In this study, we used a three-state cellular automaton market to
determine how a persistent change in one local response produced
collective selling and how dealer capacity shaped the resulting price
decline. We followed paired trajectories from a common market history
and changed one rule entry at a selected fraction of interior sites,
and we repeated the design under a held boundary that kept perimeter
traders inactive, which removed price feedback while preserving order
execution and financial accounts. Persistent activation at 20\% or 25\%
of sites produced a Sell majority in all 50 seeds within 3,000 steps
under an unlimited dealer with price feedback, whereas activation at
10\% or 15\% produced none (Supplementary Table~\ref{tab:paired-switch}).
A matched one-time Sell pulse did not reproduce the transition, while
several other single-entry rule changes did. The transition survived
removal of price feedback, so the instability lived in the lattice,
although its measured location depended on lattice size and observation
time (Supplementary Table~\ref{tab:finite-size-scaling}). With feedback
active, dealer capacity changed majority probability non-monotonically
near the activation window. Greater capacity increased drawdown at 20\%
activation, and separate cash and inventory interventions identified the
resources that limited execution. Similar Sell-dominated lattices
therefore accompanied median drawdowns from below 1\% to above 25\%,
depending on how much of the desired selling could execute
(Fig.~\ref{fig:reservoir}). Collective selling and a large price decline
were separate outcomes in this model, and the second required a dealer
able to keep buying.
